% ============================================================
%  三阶段实验时间表 — "夹心饼干→苦瓜"（2026-09-14 起，重排版）
%  编译: xelatex 20260902实验时间表.tex
% ============================================================
\documentclass[a4paper, 11pt]{ctexart}
\usepackage{geometry}
\geometry{left=2cm, right=2cm, top=2.5cm, bottom=2.5cm}
\usepackage{booktabs}
\usepackage{array}
\usepackage{enumitem}
\usepackage{xcolor}
\begin{document}

\begin{center}
{\LARGE\bfseries 三阶段实验时间表}\\[4pt]
{\normalsize ``夹心饼干 $\rightarrow$ 苦瓜''结构改造 ｜ \textbf{重排版：2026-09-14（周一）起，每天一个 mat、跨夜连续纺（$\approx$20\,h/mat）}}
\end{center}

\section*{配方速览（外层不动，只动中层）}
\begin{itemize}[itemsep=1pt, parsep=0pt]
  \item \textbf{基液（先预水解）}：三甲氧基硅烷:二甲氧基硅烷 = \textbf{5:1}（一批 75\,g = 62.5\,g + 12.5\,g；D3 支线 3:1 = 11.25+3.75）取 10\,g，\textbf{加去离子水几滴（约 0.1--0.2\,mL，不加酸）}搅拌\textbf{数小时至过夜}（水被硅烷吸收成 Si--OH）$\rightarrow$ 再与 DMF 10\,g 混匀（1:1）共 20\,g，密封防潮存放；
  \item \textbf{外层液}：基液 10\,g + \textbf{PVP 1.72\,g}（=现有外层配方；一层 mat 两外层按 20\,g + 3.44\,g），搅拌 6\,h；
  \item \textbf{中层液}：基液 10\,g + PVP 1.72\,g + 正丙醇锆 0.64\,g（$\approx$Zr:Si 5 mol\%；C-Zr$\uparrow$ 用 2.56\,g）+ 葡萄糖 0/2/4\,g，总搅拌 6\,h；
  \item \textbf{糖水}：2\,g+3\,mL、4\,g+4--5\,mL 去离子水，40--50\,°C 搅 5--10\,min 溶清，现配现用；
  \item \textbf{搅拌节奏}：硅烷+水预水解数小时$\sim$过夜；加 DMF 成基液搅 30\,min；工作液加 PVP 后共搅 6\,h；
  \item \textbf{加料顺序}：基液+PVP 搅匀 $\rightarrow$ \textbf{先加正丙醇锆}（与预水解 Si--OH 成 Si--O--Zr）$\rightarrow$ 最后加糖水，\textbf{加完尽快纺}（预水解用纯水不加酸，避免残留酸日后催化锆/糖水解）。
\end{itemize}

\section*{时间表}
\renewcommand{\arraystretch}{1.35}
\begin{tabular}{p{2.1cm}p{6.4cm}p{6.4cm}}
\toprule
日期 & 上午 & 下午/晚上 \\
\midrule
09-14（一）& 检查+小样试纺（射流稳定为准）；A/A0 中层补滴 Zr 0.64\,g（各一）；09:00 现配 A 糖水 $\rightarrow$ \textbf{09:30 开纺 A}（外--中--外，每层 10\,mL，1.5\,mL/h，约 20\,h）& 傍晚配 B、B0 中层液（各 10+1.72）搅 6\,h，21:00 各加 Zr \\
09-15（二）& $\sim$05:30 A 纺完（揭膜称初重晾置）；\textbf{A 装炉} Ar 900\,°C（跨夜）；09:00 \textbf{开纺 A0}（无糖，约 20\,h）& \\
09-16（三）& A 出炉、称重送 \textbf{SEM}；A0 纺完晾置 $\rightarrow$ \textbf{A0 装炉}；配 B/B0 外层烧杯（40+6.88）搅 6\,h & \textbf{20:00} B 糖水现配 $\rightarrow$ \textbf{开纺 B}（约 20\,h）\\
09-17（四）& A0 出炉；$\sim$16:00 B 纺完 $\rightarrow$ 预氧化 1\,h $\rightarrow$ \textbf{B 装炉} & \textbf{20:00 开纺 B0}（约 20\,h）；晚配第 2 批硅烷过夜 \\
09-18（五）& B 出炉 $\rightarrow$ \textbf{节点 1}（SEM A vs B）；第 2 批 +DMF 成基液；B0 纺完 $\rightarrow$ 预氧化 $\rightarrow$ \textbf{B0 装炉} & 配 C/E 外层烧杯与中层（6\,h），15:30 各加 Zr \\
09-19（六）& B0 出炉；09:00 C、E 糖水现配（各 4+4.5）$\rightarrow$ \textbf{09:30 开纺 C}（约 20\,h）& 晚配 Zr 系列外层烧杯与 C-Zr$\downarrow$/C-Zr$\uparrow$ 中层 \\
09-20（日）& $\sim$05:30 C 纺完 $\rightarrow$ 预氧化 $\rightarrow$ \textbf{C 装炉}；09:00 \textbf{开纺 E}（不预氧）& \\
09-21（一）& C 出炉 $\rightarrow$ SEM+BET \textbf{节点 2}；E 纺完 $\rightarrow$ \textbf{E 装炉}；09:00 \textbf{开纺 C-Zr$\downarrow$}（糖水现配）& \\
09-22（二）& E 出炉；C-Zr$\downarrow$ 纺完 $\rightarrow$ 预氧化 $\rightarrow$ \textbf{C-Zr$\downarrow$ 装炉}；09:00 \textbf{开纺 C-Zr$\uparrow$} & 晚配 D3 硅烷小批（3:1）过夜 \\
09-23（三）& C-Zr$\downarrow$ 出炉；C-Zr$\uparrow$ 纺完 $\rightarrow$ 预氧化 $\rightarrow$ \textbf{C-Zr$\uparrow$ 装炉}；D3 加 DMF 成基液 & 配 D3 工作液（6\,h）\\
09-24（四）& C-Zr$\uparrow$ 出炉 $\rightarrow$ Raman + \textbf{节点 3}；09:00 \textbf{开纺 D3}（约 20\,h）& \\
09-25（五）& D3 纺完 $\rightarrow$ 预氧化 60\,min $\rightarrow$ \textbf{D3 装炉} & 缓冲；可补 B-t0.5 / B-t2 \\
09-26（六）& D3 出炉；协调 1250\,°C 炉（C/E 各切一份）& 集中表征准备 \\
09-27（日）& 送 \textbf{C-1250 / E-1250} 至原 1250\,°C 炉 & \\
09-28（一）& 1250 支线出炉 & VNA 介电（算 $|Z_{in}/Z_0|$ 与 $\alpha$）+ 热导率 \\
09-29（二）& 补测复核 & 数据整理 \\
09-30（三）& \textbf{汇总}：形貌--阻抗--隔热/吸波 联表，写小结 & \\
\bottomrule
\end{tabular}

\section*{D3 支线（B 配方、硅烷 3:1）穿插执行}
\begin{itemize}[itemsep=1pt, parsep=0pt]
  \item 09-22（二）晚：配 3:1 硅烷小批 15\,g（MTMS 11.25\,g + DMDMS 3.75\,g）+ 几滴水预水解过夜；
  \item 09-23：加 DMF 15\,g 成基液；09-23 下午配 B 配方工作液（外层：基液 20\,g + PVP 3.44\,g；中层：基液 10\,g + PVP 1.72\,g + Zr 0.64\,g + 糖 2\,g 糖水），搅拌 6\,h；
  \item 09-24 09:00 现配糖水后开纺 D3（外--中--外）$\rightarrow$ 09-25 纺完预氧化 220\,°C 1\,h $\rightarrow$ 装炉 Ar 900\,°C $\rightarrow$ 09-26 出炉；
  \item 目的：与 B 组（5:1）同配方直接对比可纺性与 ID/IG，定下一轮是否整体切 3:1。
\end{itemize}

\section*{三个决策节点}
\begin{enumerate}[itemsep=2pt]
  \item \textbf{09-18}：预氧化是否成壳/表面出现差异？否 $\rightarrow$ 调预氧化温度或时间（B-t2）再验；
  \item \textbf{09-21}：糖量×预氧化对孔隙（BET）与形貌影响？Raman 看自由碳是否被糖抬升；
  \item \textbf{09-24}：Zr 掺量对碳保留（ID/IG）的影响，定保碳-耐温平衡点。
\end{enumerate}

\section*{贯穿红线}
\begin{itemize}[itemsep=1pt, parsep=0pt]
  \item \textbf{纺丝节奏（按历史速率）}：1.5\,mL/h 下每层 10\,mL $\approx$ 6.7\,h，一个 mat（外--中--外）$\approx$ \textbf{20\,h 连续纺} $\rightarrow$ \textbf{每天只起一个 mat}，需\textbf{跨夜无人值守}（电压/流速稳定、硅油纸够长、针头不堵）；不能无人值守就把每层降到 3--4\,mL；
  \item \textbf{试纺判定}：以\textbf{射流稳定}为准（泰勒锥稳、无堵针/滴液、射流连续），约 5--15\,min（$\approx$0.1--0.4\,mL 弃去）即可开收，不按体积算；
  \item \textbf{停置检查（必做）}：基液与 PVP 液自 09-04 密封至今（9 天），纺前逐个烧杯目视——透明无浑浊、无凝胶块、倾倒无拉丝；\textbf{先取 2--3\,mL 小样试纺}，增稠/凝胶则\textbf{重配基液}，不要硬纺；
  \item \textbf{烧杯规范}：搅拌与存放用保鲜膜/封口膜封口（防 DMF 挥发吸湿），杯身贴标签并与实验日志「烧杯记号」栏一致；纺前每杯做一次定性黏度记录（挑丝/挂壁/气泡悬浮）；
  \item \textbf{控水}：基液密封；葡萄糖烘干现用；正丙醇锆在糖水前加入（与 Si--OH 结合）、防潮；加完糖水尽快纺；记录环境状况（有湿度计则记数值）；
  \item 纺丝液配好先做一次定性记录：玻璃棒挑丝长度、挂壁厚薄、气泡悬浮时间、倾倒流速（本组无电导率仪，用目视法代替电导率对照用水/缩合问题）；
  \item 预氧化 $\le$ 250\,°C，别烧掉中层糖；空气与 Ar 两段气氛分开；
  \item 每组记录实际葡萄糖/水/正丙醇锆质量、初重与失重；
  \item 无糖对照（A0/B0）必须做；阻抗匹配判定先行，再谈 RL。
\end{itemize}

\end{document}
